% id: 117-59
% book: 쎈B 공통수학2 · 08 함수
% section: 유형 14 $f^n$ 꼴의 합성함수
% figure: none
% answer: 
% answer_source: 
% variant_level: 0 (원본 전사)
% 이 파일은 build.mjs 가 items.json 에서 만든다. 고칠 때는 items.json 을 고친다.
\smash{\raisebox{6.5mm}{\dmkichul{쎈B 공통수학2 117쪽 59번}}}
자연수 전체의 집합에서 정의된 함수 $f(x)$가\par\smallskip\dispeq{$f(x)=\begin{cases} \dfrac{x+2}{3} & (x=3k-2) \\[1mm] \dfrac{x+1}{3} & (x=3k-1) \\[1mm] \dfrac{x}{3} & (x=3k) \end{cases}$\ ($k$는 자연수)}\par\smallskip\noindent 이고 $f^1=f$, $f^{n+1}=f^n\circ f$로 정의할 때, $f^n(99)=1$을 만족시키는 자연수~$n$의 최솟값은?
\dmchoicesauto{}{$3$}{$4$}{$5$}{$6$}{$7$}
